\quellseite{930}
\chapter[C. D. Mart. Luthers Auslegung über das Lied Mosis,]{C. D. Mart. Luthers Auslegung über das Lied Mosis, 5 Mose 32.\footnote{Luther hat diese Auslegung ursprünglich lateinisch verfaßt, doch ist von derselben weder eine lateinische Einzelausgabe bekannt, noch ist sie in die lateinischen Sammlungen der Wittenberger und der Jenaer Ausgabe aufgenommen. Wann die Auslegung geschehen sei, ist auch nicht genau bekannt. Seckendorf (Hist. Luth. lib. II, § 53, p. 148a) und Walch (Vorrede zum 3. Bande, § 12) setzen sie in das Jahr 1529; doch kann es kaum einem Zweifel unterliegen, daß dieselbe später zu setzen ist. (Vgl. die Anmerkungen zu § 29, § 50 und § 70.) Justus Jonas übersetzte dieselbe ins Deutsche und versah sie mit einer Vorrede. Diese Uebersetzung (vielleicht nach einem Manuscripte) erschien zuerst bei Georg Rhaw zu Wittenberg im Jahre 1532 unter dem Titel: Auslegung D. Mart. Luthers über das Lied Mose am zwei und dreißigsten Cap. Deutero. Verdeutscht aus dem Latein durch Justum Jonam. Noch in demselben Jahre kam eine andere Ausgabe heraus, gedruckt zu Nürnberg durch Kunigund Heragotin. In den deutschen Sammlungen findet sich diese Uebersetzung: in der Wittenberger (1556), Bd. V, Bl. 247 b; in der Eislebenschen, Bd. I, Bl. 555; in der Altenburger, Bd. IV, S. 786; in der Leipziger, Bd. IV, S. 287, und in der Erlanger, Bd. 52, S. 400. Georg Walch als auch die Erlanger Ausgabe haben den Text der ersten Ausgabe zu Grunde gelegt, daher sind wir diesen gefolgt. Die Wittenberger Ausgabe hat sich so viele und zum Theil so bedeutende Veränderungen mit dem Texte erlaubt, daß sie für uns nur hie und da zur Richtigstellung einzelner Lesarten in Betracht kommen kann. Wenn man nur die Wittenberger Ausgabe hätte, würde man glauben, daß auch die Vorrede von Luther herrühre, denn die Ueberschrift derselben lautet: „Vorrede in den Gesang Mose“, und in der Vorrede selbst sind die Worte: „nach der Auslegung D. Martini“ weggelassen worden, und der betreffende Passus lautet daselbst: „hab ich dieses Lied derhalben deutsch in Druck gegeben“. In der Erlanger fehlt die Vorrede des Jonas.}}

\subsection{Ausgelegt 1530 oder später, ausgegangen 1532}

Aus dem Lateinischen übersetzt durch Justus Jonas.

\section{Vorrede des Justus Jonas}

Als der hohe Prophet Mose in seinem fünften Buche die zehn Gebote, nämlich der ersten und andern Tafel, durch viel edle, treue, köstliche Predigten dem Volke hatte eingebildet und ausgelegt; wie denn dasselbige fünfte Buch ein rechter Ausbund ist, und eine überaus vortreffliche glossa ordinaria und commentarius über die zehn Gebote; hat er kurz vor seinem Tode dem ganzen Volke Israel und allen Nachkommen dieses Lied zu einer Gedächtnis und zur Lehre gelassen, darinnen er mit kurzen Worten die Summa des ersten, vornehmsten und höchsten Gebots, und sonst aller Predigten, Vermahnung und Tröstung, so er lange Zeit, vom ersten Tage seines Berufs und Eingangs, bei dem großen Volke gethan, gefasset und begriffen hat; damit sie solches erstes, höchstes Gebot der ersten Tafel: „Ich bin dein Gott“ etc. und das ganze göttliche Gesetz, nicht allein von Priestern und Leviten alle sieben Jahre, und zu seiner Zeit hörten ausrufen, vorlesen und predigen, sondern daß auch dieselbige Lehre unter dem Volke, bei Jungen und Alten, in der Gemeinde und in Häusern, mit täglicher Erinnerung, mit täglichem Lehren, Leben und Vermahnen, mit täglichem Singen und Sagen getrieben und den Leuten eingebläuet würde; wie man jetzund noch Lieder und Gesänge dem gemeinen Manne und der Jugend macht, darinnen die Summa des ganzen Evangelii und Erkenntniß Christi begriffen ist. Denn Mose, als ein besonderer hoher Mann, hat wohl gewußt, daß diese selige Lehre, nämlich das rechte GOttes Wort, im Volke auszubreiten und zu erhalten, keine menschliche Weisheit, keine Arbeit, keine Sorge, Mühe, noch Fleiß genug ist, und daß es auch nicht hilft, daß man die Bücher und Schriften hat; wie etliche jetzund wähnen, wenn sie allein die Bücher haben, so sind sie gelehrt und heilig; ja, daß es nicht hilft, daß das Deuteronomion in der Lade des Bundes, oder in der Liberei liege [2 Mos. 25, 16. 40, 20. 2 Kön. 22, 8.], sondern wenn auch die Prediger (als solcher hoher Prophet aller hohen Propheten, Moses selber, und andere) das göttliche Wort mit allem höchsten treuen Fleiße, mit lebendiger Stimme und täglichem Anhalten treiben. Wenn diese zehn Gebote oder göttlich Wort gleich täglich singen und alten Leuten mit täglichem Anregen, mit Bitten, mit Flehen; mit Loben, mit Singen und Sagen, mit aller innerlicher und äußerlicher Uebung, werden eingebildet, darf es dennoch großes Glück, und sonderbarer reicher Gnaden GOttes, daß es vor so großem, grimmigem Toben und Wüthen des Teufels, der es gerne zu Grunde tilgen wollte, doch bei etlichen hafte und bleibe, und durch falsche
\quellseite{931}
Lehre und andere List des Satans und Bosheit der verzweifelten Welt nicht ganz vergessen und gar unterdrückt werde.

Und so nun Moses, als ein hoher Mann, welchem so eine unsäglich große Bürde, das große Volk zu regieren, war aufgelegt, kurz vor seinem Absterben schon in seinem Herzen mit den Gedanken umging, wie bald alle das Gute, das er angerichtet hatte, würde fallen; wie jämmerlich und erbärmlich es nach seinem Tode und weiter hernach würde zugehen; wie es je länger je ärger in demselbigen Volk künftig werden sollte, hat er die große, treffliche Wohlthat, so GOtt dem Volke erzeiget, in diesem Liede kurz mit ganz wichtigen Worten errührt und erholet;\footnote{erholen = wiederholen.} darnach durch das ganze Lied bis zum Ende setzt er, als ein rechter, scharfer Moses, eitel harte, schreckliche Strafworte und greuliche Träumung göttliches Ernstes und Zorns.

Und, wie das seine Sprüchwort lautet [Matth. 12, 34.]: „Wes das Herz voll ist, gehet der Mund über“, sieht man wohl in diesen Worten, daß Moses dieses Lied aus heftigem, hitzigem, göttlichem Eifer geschrieben hat, und daß dies Lied kurz vor seinem Tode das treue fromme Herz, den trefflichen hohen Mann, manchen schweren, betrübten Gedanken, manche starke Anfechtung gekostet hat; wie wir denn sehen werden. Als, wenn jemand ein treuer frommer Prediger, der aus dem Grunde seines Herzens gerne sähe, daß GOttes Wort unter dem Volke bleiben möchte und auf die Nachkommen reichen, ein Lieblein säen sollte, wie es noch in dreißig oder vierzig Jahren wird in aller Welt zugehen, um der trefflichen, unsäglichen, großen Undankbarkeit und Verachtung willen, welche jetzund Bischöfe und Fürsten, und auch Edelleute, Bürger, Bauern etc. auf dieser Seite gegen das göttliche Wort erzeigen, der würde auch nicht viel Freude oder Gaudeamus, sondern viel Herzleid, Requiem und Jammer zu singen finden. Vielmehr hat Moses große Angst und Schmerzen davon gehabt, welcher hier in diesem Liede prophezeiet, wie es dem ganzen Volke künftig und über lange Zeit, wenn sie nun ihre Könige und Fürsten haben werden etc., gehen soll, bis daß [es] gar mit ihnen aus sei.

Dieweil aber dies Lied nicht allein die Israeliten angeht, sondern auch uns alle, und die lieben zarten Gesellen, welche jetzund in allen Ständen GOtt, den HErrn, mit seinen zehn Geboten verachten, und sagen: Die Pfarrherren können nichts anders, denn die zehn Gebote predigen, habe ich dies Lied, nach der Auslegung D. Martini,\footnote{Die Worte: „nach ... Martini“ fehlen in der Wittenberger Ausgabe.} deutsch wollen an Tag geben, die Gelehrten, so ihre Pfarrherren mit den zehn Geboten verachten, wieder umgekehrt zu machen; anzuzeigen, daß Moses die zehn Gebote nicht geringe achtet, daß sie doch sehen, was die zehn Gebote sind, wer der Mann ist, der auf dem Berge Sinai (wie 2 Mos. 19, 16. 18. beschrieben) mit großem schrecklichen Feuer, mit starkem Blitz und Donner, als einer, der HErr sein und bleiben will, sich hat sehen und hören lassen; und was Großes der ganzen Welt, Kaiser, Papst, Fürsten, Herren, Königreichen, Landen, Leuten an den kurzen zehn Geboten (welche ein jeglicher Bauer und Bauerknecht jetzund hören kann, aber gleichwohl verachtet, und mit Füßen tritt) gelegen ist.

Es ist eine große, sonderliche GOttes Gabe und Gnade, daß jetzund wiederum der Katechismus rein gelehret wird in der Kirche, nämlich, die zehn Gebote, Glaube, und Vater-Unser. Denn dadurch hat ein jeglicher Mensch die ganze Theologie, und kann erkennen lernen, was der rechte, beste, höchste, wahrhaftige Gottesdienst ist, was GOttes Wille und Gebot ist, was er von uns fordert; item, wie ein jeder in seinem Stande, in allem Wandel und Leben, gegen seinem Nächsten sich halten soll, und leben, daß es GOtt gefalle. Und wenn die lutherische Lehre (welcher der Teufel so bitterlich feind ist) nichts anders genützt hätte, denn daß sie den Katechismum und die zehn Gebote hat wider dem Volk bekannt gemacht, welcher unter dem Papsthume gegen den Wallfahrten, gegen der Lehre von Möncherei und Messe hören etc. mußte die geringste Lehre sein, so hätte sie doch mehr in der christlichen Kirche gebauet, denn Paris und alle hohe Schulen, so lange sie auf Erden gewesen.

Aber die schändliche Welt ist allenthalben undankbar und verkehrt, beide gegen GOttes Gebote und gegen dem Evangelio. Hört sie Gebote, so kann sie nicht denken, wie groß der Gebieter ist, bis daß sie Blitz und Donner fühlt, entweder leiblich, oder im Gewissen geistlich. Verkündigt man ihr denn Gnade und Barmherzigkeit, so kann sie aber\footnote{aber = abermals.} ihr nicht schenken lassen, sondern will mit ihren eigenen Werken GOtt Gnade abverdienen.

Also, wenn man den Leuten die zehn Gebote predigt, denken sie nicht, daß es GOtt mit den Geboten so ein großer Ernst ist, und daß sie alle die größten Sachen, die im Himmel und Erden, in dieser und jener Welt sind, belangen. Denn die erste Tafel belangt je GOtt selber, seine Gottheit, seine göttliche Ehre, welche ihm Engel und Menschen schuldig sind. Die andere Tafel belangt aber die größten Sachen, die unter dem Himmel und auf Erden sind, alle Politien und Regimenter, alle
\quellseite{932}
Haushaltung, alle Händel und Werbung, und ist also keine Weisheit, kein Leben oder Wesen, das nicht in die Gebote gefasset sei.

Viel aber, wenn sie die zehn Gebote hören nennen, denken an der Kinder Bibel, meinen, es sei eine schlechte Kinderpredigt; denken nicht, daß es so ganz hohe, große Gebote sind, und daß hinter denen so ein großer HErr, so eine hohe Majestät, so ein hoher Gebieter ist. Aber wir wissen, GOtt Lob, daß die zehn Gebote hören und lernen, item, wissen, was Gesetz, was Evangelium, was Gebot, was Verheißung ist, nicht so eine schlechte Kinderlehre zu achten ist, sondern es sind die höchsten zwei Lehren unter der Sonne, ohne welche niemand zu GOtt kommen, oder die Wahrheit erkennen mag.

Von dem Evangelio zu reden wollen wir auf diesmal lassen anstehen; aber wenn du die zehn Gebote hörst nennen, darfst du nicht mit deinem Gedanken in der Kinder Bibel fahren, oder ansehen, daß es kurze Worte sind, welche in andern Kinderbüchern auch geschrieben stehen, welche man täglich hört, sondern magst wohl höher fahren, und wissen, daß du eine Lehre hörst, sonderlich im ersten Gebote: „Ich bin dein Gott“ \&c., „du sollst nicht andere Götter haben“, die noch kein Prophet, auch Moses selber, kein Apostel, ja, kein Engel im Himmel ausgelernt oder ausstudirt hat.

Denn, daß eine arme Creatur, ein armer Mensch, die hohe Majestät, welche Himmel und Erde, Engel und Menschen geschaffen hat, wahrlich für ihren HErrn und GOtt halte, sich alles Gutes zu ihm versehe, sich auf ihn wage im Leben und Sterben, ist nicht eine Kunst, die man kann ausstudiren. So hat auch kein Heiliger auf Erden je gelebt, auch nicht Johannes Baptista, der sie ausstudirt habe. Darum sagt auch der HErr Christus Matth. 22, 40., daß in den zweien Geboten das ganze Gesetz hange, und alle Propheten.

Denn im ersten Gebot ist alles begriffen, was GOttes Ehre, seine Gottheit, seinen heiligen Namen, seinen heiligen Willen, und den höchsten, heiligsten Gottesdienst angeht. Darum alle Ceremonien, alles geistliche Leben, alle Gottesdienste, sie sind unter Christen oder Heiden, sollen und müssen sich durch das erste Gebot urtheilen und richten lassen, also daß, wo die dem gemäß sind, so sind sie GOtt recht wohlgefällig; sind sie dem Gebot ungemäß, so sind sie abgöttisch, teuflisch und eitel Greuel vor GOtt.

Also im andern Gebote: „Liebe deinen Nächsten“, oder in der andern Tafel, sind begriffen alle rechten, göttlichen, christlichen guten Werke, die in aller Welt von Jungen oder Alten, in was Stand sie sind, immer geschehen mögen, sollen auch und müssen durch die andere Tafel sich urtheilen und richten lassen. Und die viel von guten Werken schreien, aber dieses größte Register, da die guten Werke erzählt werden, mit Füßen treten, Gebruch, Mord, Dieberei ohne Unterlaß treiben, mögen sich wohl spiegeln in der andern Tafel. Denn alle anderen guten Werke, ob sie wohl groß scheinen und gleichen, wenn sie den Geboten der andern Tafel nicht gemäß sind und im Gehorsam GOttes nach dem ersten Gebot geschehen, sind eitel heuchlerische Werke, derer GOtt nicht achtet, und die dem Nächsten nicht nütze sind.

So ist es nun alles, was GOtt und den Nächsten belangt, in die zehn Gebote gefasset; und aller Propheten Bücher, auch viel Psalmen, sind eitel Predigt, Erinnerung und Auslegung, mit mehrern und reichern Worten, des ersten Gebots: „Ich bin dein Gott“ \&c. Item, alle Historien im Josua, im Buch der Richter, der Könige, der Chroniken sind Exempel, wie GOtt so fest über dem ersten Gebote gehalten hat, wie er die Ungläubigen und Gottlosen gestraft, und wie er die Gläubigen errettet und getröstet hat.

Aber also groß ist die List und Bosheit des Teufels, daß er die Heiligkeit und Gottesdienst, welche vor der Welt den größten Ruhm und Schein haben, als, die unzähligen Satzungen und Mönchereien unter dem Papst und Türken, kann anrichten, daß jedermann, eher das mehrere Theil davon fällt, und die rechten Gebote GOttes, sonderlich das erste, nicht allein läßt anstehen, sondern mit Füßen tritt.

Also ist unter dem Papstthum kein toller Mönchstraum gewesen, keine Lüge von Rosenkränzen so unverschämt, sie ist heftiger mit Lehren und Predigen, mit Worten und Werken getrieben, denn die heiligen zehn Gebote, und das höchste, heiligste, vornehmste Gebot, welches den höchsten, reinsten, vornehmsten Gottesdienst fordert, daß wir GOtt von Herzen glauben sollen, ihn fürchten \&c.

Darum, so nun Moses wohl wußte, daß im Volk Israel auch allerlei Heuchler würden aufstehen, dem Gottesdienste Baal, und andern falschen Gottesdiensten anhangen, von diesen Geboten abweichen, und das erste Gebot in den Wind schlagen, so hat er ihnen dieses Lied zur Letze gelassen, darinnen er ihnen mit kurzen, aber ganz heftigen Worten dräuet, daß ihnen hernach alle Propheten gedräuet haben, welches ihnen auch also widerfahren ist, daß beide Himmel und Erde zeugen müssen, es habe kein Tüttel gefehlt, es sei alles wahr worden.

Und in diesem Liede werden wir nun sehen, daß GOtt mit den zehn Geboten keinen Scherz will haben, sondern, wie gering und schwach sie scheinen, wie wenig sie von Gewaltigen, von Reichen, von Weisen, von den falschen Heiligen geachtet sind, so sollen sie doch im Himmel und auf Erden regieren
\quellseite{933}
, alle anderen Gebote und Lehren, so den Weg der Seligkeit wollen weisen, alle falsche Heiligkeit und Ketzereien verdammen, Königreiche und Fürstenthümer umstoßen und aufrichten; wie sie denn am Königreiche Israel und Juda bewiesen haben, und noch täglich beweisen mit vielen Exempeln.

Denn, ob es die Welt nicht glaubt, und in Sicherheit und Blindheit GOtt mit seinen Geboten verachtet, so fährt GOtt, der HErr, gleichwohl fort, er zeigt, daß er der HErr und GOtt ist, tröstet und errettet seine frommen Kinder, erzeigt ihnen reiche Güte und Gnade; wiewohl unter allerlei Kreuz und Trübsal. Wiederum plagt, straft, tilgt, und stürzt er endlich die Gottlosen, wie das erste Gebot lautet, und wie die Psalmen davon singen; wie in allen Historien, beide der Schrift und der Heiden, zu sehen ist. Das werden wir in diesem Liede weiter hören und lesen.

\section{Das Lied Mosis, 5 Mose 32}

1. In diesem Gesange oder Liede (wiewohl er zur Lehre billig freundlich reden sollte) braucht der Prophet Mose nicht viel sanfter, gelinder Worte, sondern redet alles daher mit eitel großem, trefflichem, hartem Ernst; schilt, dräuet und straft aufs geschwindeste, erzeigt eitel Härtigkeit und Schärfe. Und erstlich, wirft er dem undankbaren Volke vor GOttes unzählige, großmächtige, leibliche und geistliche Wohlthaten, stellt ihnen dieselben vor die Augen, daß sie es greifen mögen. Wiederum malt er ihnen auch vor die Augen ihre feine Dankbarkeit, welche sie und ihre Nachkommen dem frommen, treuen GOtt beweisen werden; nämlich, daß sie aller großen Wohlthaten, und GOttes selber vergessen werden, sich abkehren zu fremden Göttern, und derhalben greulichen GOttes Zorn und Vermaledeyung auf sich laden.

2. Und Moses thut hier wie ein rechter Moses und Gesetzprediger, daß er nicht viel väterlich lockt, sondern schilt, straft, eifert und zürnt über ihre große Undankbarkeit, redet eitel brennende, feurige Worte, dräuet allenthalben mit einem allmächtigen, starken GOtt, deß Händen niemand entgehen noch entfliehen werde.

3. Damit zeigt er an die Art des Gesetzes [Röm. 4, 15. 7. 7.], nämlich, daß es die Sünde zeigt, das Gewissen erschreckt, und daß ohne Gnade und Barmherzigkeit durchs Gesetz niemand vor GOtt kann fromm werden, daß [es] auch nicht zu wagen sei, mit einer solchen großen Majestät durch unsere bettlichen, elenden Werke zu handeln, oder dem Feuer mit Strohhalmen zu begegnen.

\begin{schriftzitat}
V. 1. Merket auf, ihr Himmel, ich will reden, und die Erde höre die Rede meines Mundes.
\end{schriftzitat}

4. Er ruft Himmel und Erde und alle Creaturen an zu Zeugen, als in einer großen Sache, welche nicht zeitliche Güter, nicht ein Königreich oder Fürstenthum, nicht ein Haus, Hof, Gold oder Silber, sondern einem jeglichen das Höchste und Beste belangt; nämlich GOtt, den Brunn alles Guten, mit allem ewigen und zeitlichen Heil, mit Himmel und Erde und allen Creaturen, entweder zu haben, oder ewig (welches schrecklich ist) zu verlieren und zu entbehren. Ja, welche Sache nicht belangt den Kaiser, nicht Fürsten oder Herren zu versöhnen, sondern GOtt, die ewige Majestät, gnädig zu haben, hier und dort wohl zu fahren, oder in GOttes Ungnaden, in ewigem Ver zagen an allem göttlichen und menschlichen Troste, ohne alle Hülfe oder Rettung, ewig und ohne Unterlaß zu sterben und verderben.

5. In solcher großen Sache sollte noch wohl klein Zeugniß oder kleine Zeugen zu wenig sein; da müßten noch wohl vertreten der schöne, hohe, weite Himmel, mit der edlen Sonne, mit dem Mond, mit allen Sternen. Da müßte noch wohl die Erde mit allen ihren Gewächsen, mit allen Vögeln, mit allen Thieren kommen, und das große weite Meer mit alle seinen Fischen, und alles, was sich darinne reget, und ihrem GOtt wider die Gottlosen Zeugniß geben, seine ewige göttliche Ehre und Gerechtigkeit erretten, und sein Urtheil bestätigen [Offenb. 6, 12. Weish. 19, 6. ff.].

6. Das geht auch also: Wenn das Stündlein kommt, daß GOtt, der HErr, wider die Verächter aufwacht, und sich nimmer will vergessen oder verachten lassen, entfällt ihnen plötzlich alle ihr Muth, werden in ihrem Herzen und Gemüssen
\quellseite{934}
so ganz verstürzt, blöde, traurig und verzagt, daß sie nichts anderes fühlen, denn daß nicht allein GOtt, sondern Himmel, Erde, Laub und Gras, alle Creaturen mit ihnen zürnen, sie überzeugen, anklagen und verdammen.

7. Also dräuet der Prophet Habakuk dem Könige zu Babylon, da er sagt am 2. Capitel, V. 11.: „Denn auch die Steine in der Mauer werden schreien, und die Balken am Gesperre werden ihm antworten“; das ist, wenn GOtt ihn werde heimsuchen, solle er so einen verzagten Muth kriegen, wenn ein Brett an der Wand knacket, daß er denke, der Himmel wolle fallen. Also dräuet auch Mose am 26. des dritten Buchs, V. 36., den Gottlosen, daß sie GOtt werde so verzagt machen, daß sie soll ein rauschend Blatt jagen etc. Dergleichen aber sehen wir Jesaia am 1. Capitel, V. 2., daß er, dem Mose nach, über die Undankbarkeit des halsstarrigen, gottlosen, bösen Volks Himmel und Erde anruft. Ja, sie werden auch allzu hart zu seiner Zeit zeugen.

\begin{schriftzitat}
V. 2. Meine Lehre triefe wie der Regen, und meine Rede wie der Thau, wie der Regen auf das Gras, und wie die Tropfen auf das Kraut.
\end{schriftzitat}

8. Es wünscht hie Moses, daß seine Lehre möge triefen, wie der Regen und Thau auf das Gras und Kraut, das ist, daß sie möge kräftig sein und Frucht bringen, daß er nicht in Wind rede und predige. Nachdem er von der allergrößten Sache auf Erden redet, vom Handel der Religion, nämlich vom ersten, höchsten Gebote GOttes, was die rechte Religion, was der höchste und trefflichste Gottesdienst ist. Und wiederum, was der höchste Greuel vor GOtt sei; nämlich falsche Lehre, falscher Gottesdienst, falsche Heiligkeit, allerlei Heuchelei, welches alles stracks wider das erste, höchste Gebot GOttes geht, welche Heuchelei Moses [V. 33.] nennt eitel Ottern und Drachen Gift und Galle [Jes. 59, 5. 6.].

9. So wünscht er nun, daß diese seine rechte Lehre bei vielen möge Frucht schaffen. GOttes Wort vergleicht er [dem] Regen und Thau. Denn wo es rein gepredigt wird, da ist es nicht, wie menschliche Satzung, ohne Nutz oder Frucht, welche Petrus [2. Ep. 2, 17.] „Wolken ohne Regen“ nennt, sondern wirkt und schafft etwas Großes. Die Zuhörer vergleicht er dem Gras und Kraut [Jes. 40, 6.]. Denn wie Regen und Thau die Wiesen lustig und grün machen, daß Blumen und Gras nach aller Lust, nach allem Wunsch aufgehen und wachsen; also erquickt GOttes Wort die Herzen und Gewissen; und wo derselbige Regen göttliches Worts hinfällt, da geht es ohne Besserung und Frucht nicht ab.

\begin{schriftzitat}
V. 3. Denn ich will den Namen des HErrn anrufen. Lasset unsern Gott die Herrlichkeit.
\end{schriftzitat}

10. Das ist, ich will ein Liedlein singen, das will ich hoch anheben, und höher soll es niemand unter der Sonne anfangen, noch besser machen können. Mein bestes Lied und beste Lehre soll sein das erste Gebot GOttes; nämlich, daß man dem einigen, wahren GOtt die Ehre geben soll, ihn herzlich fürchten und lieben, auf ihn allein trauen und bauen, und daß alle anderen Gottesdienste, wie heilig, fröhlich und köstlich sie scheinen, so außer oder neben dem ersten Gebot werden angerichtet, eitel Abgötterei und Heuchelei sind.

11. Und lautet mein Lied also: „Gebet GOtt die Herrlichkeit.“ Das ist, dem einigen, wahren, lebendigen GOtt gebührt allein die Herrlichkeit, Lob und Preis im Himmel und auf Erden. „Er ist allein (wie Moses sagt in diesem fünften Buch Cap. 10, 17.) ein GOtt über alle Götter, ein HErr aller Herren“, der Himmel und Erde und das Meer geschaffen hat, und alles, was darinnen ist [1 Tim. 6, 16. Offenb. 1, 8.]. Der alle Königreiche auf Erden in der Hand hat, befehliget, ausrichtet, und niedersetzt, wie er will; der allen Menschen auf Erden Leben und Odem gibt; der aller Könige Herzen, Sinne und Gedanken, und sonst aller Menschen auf Erden, inwendig bildet und lenkt, wie es ihm gefällt; der allein alle leiblichen und geistlichen Gaben gibt; ohne welchen niemand Leib, Leben, Weisheit, Stärke, Gesundheit, Gewalt, Reichthum oder einiges Gut haben, oder einen Augenblick behalten kann. Denn, spricht Moses, gebet die Herrlichkeit, das ist, ihn erkennet für einen HErrn, der alles schafft und wirkt, von dem allein alle Gaben herrühren. Also sagt auch Jeremias am 9., V. 23. 24.: „Ein Weiser rühme sich nicht seiner Weisheit, ein Starker rühme sich nicht seiner Stärke, ein Reicher rühme sich nicht seines Reichthums, sondern wer sich rühmen will, der rühme sich des, daß er mich wisse und kenne, daß ich der Herr bin“ etc.

\quellseite{935}
12. Hier ist nun beschlossen, daß keines Menschen die Herrlichkeit ist, sondern alle Hoheit, alle Majestät, Gewalt und Macht ist allein GOttes, der gibt allein Weisheit [Luc. 21, 15.], Stärke, Reichthum, wenn er will, nimmt es auch wiederum weg, und ändert es mit allen Menschen, wie er will. Darum kann kein Königreich länger stehen, denn er will; kein Regiment, keine Haushaltung, kein Reichthum oder Wohlfahrt länger währen oder bleiben, denn er will. Wie denn Daniel Cap. 2, 21. sagt: Er richtet Königreiche auf, und stößt sie zu Boden, wie er will.

13. Das nun also erkennen und wissen, ihm allein alle Ehre heimweisen, ernstlich von ihm alle Gaben, leiblich, geistlich alle Hülfe und Trost [er]warten, auf ihn in Glück und Unglück, in Leben und Sterben, sich von Herzen verlassen, zu ihm in allen Anfechtungen fliehen, ihn allein in allen Nöthen und Trübsalen suchen und anrufen, das ist der höchste, angenehmste Gottesdienst.

14. Die Herrlichkeit aber gibt\footnote{Erlanger: „gebe“.} Gott, dem HErrn, allein der kleine Haufe, die Gläubigen und Frommen auf Erden; dieselbigen vertrauen und bauen auf ihn von Herzen, erkennen, daß sie alle Güter von ihm haben, daß sie keinen Augenblick vor dem Teufel bleiben könnten, wenn GOtt sie nicht schützt und behütet [1 Sam. 2, 6–8.]. Die Heuchler aber, und alle anderen Gottlosen, ob sie wohl mit dem Munde GOttes Allmächtigkeit und Herrlichkeit bekennen, so verlassen sie sich doch im Herzen auf ihre eigene Heiligkeit, Weisheit, Stärke, Vermögen, auf ihre eigenen Werke, auf Menschen Hülfe und Trost; und wenn sie in Nöthen dieselbigen nimmer empfinden, so erschrecken sie, verzagen und verzweifeln, laufen von GOtt, lästern und fluchen Gott.

15. Dieses ist aber ein starker Trost den Frommen und Gläubigen: Ob sie wohl täglich in der Welt sehen, daß GOttes Wort und Werke so schändlich verachtet, verlachet, verachtet und mit Füßen getreten werden, daß sie dennoch für wahr wissen, daß es mag eine Zeitlang währen. Es mag einen Augenblick GOtt den Verächtern zusehen, aber endlich erfahren sie, daß dies erste Gebot wahr bleibt, und die Macht und Herrlichkeit GOtt behält. Wie man jetzund sieht, daß Klein und Groß ohne Gottesfurcht, ohne Scheu geizen, GOttes Wort und die frommen, treuen Prediger verachten, und die zehn Gebote als eine geringe Lehre in Wind schlagen; wissen aber nicht, daß endlich dies erste Gebot wird wahr bleiben; und daß sie jetzund nicht glauben wollen, das werden sie in ihrem Gewissen, an Leib und Seele, an Ehre und Gut, Weib und Kindern erfahren müssen.

16. Darum, wenn man GOtt die Herrlichkeit läßt, alle Hülfe und Trost bei ihm sucht, und von ihm wartet, da hat er Segen und Gnade zugejagt, da geht es auch wohl, es sei in Kirchensachen oder Predigtamte, es sei in Regimenten, in Haushaltung, kleinen oder großen Händeln. Das sehen wir: Wo ein frommer, treuer Bischof und Pfarrherr ist, der nicht eigene Ehre sucht, sondern erkennt, daß er den Geist, Gnade, das Wort und alle Gaben von GOtt hat, da ist GOtt selber Bischof und Pfarrherr mit; wiederum, wo Ketzer und falsche Lehrer GOtt die Ehre rauben, und ihre eigene Ehre suchen, da geht es, wie wir zu unserer Zeit an Thoma Müntzer und andern mehr gesehen und erfahren haben. Also auch in Regimenten: Wo gottesfürchtige Obrigkeit ist, die GOtt vor Augen hält, ihm die Herrlichkeit läßt, da sitzt GOtt selber mit im Rath, hilft mit zu sehen, wachen und regieren.

17. Wiederum: Wo die Könige, Fürsten und Herren GOttes vergessen, rauben GOtt seine Ehre, bauen auf ihre eigenen Anschläge und Weisheit, auf Reichthum und Gewalt, auf Bündnisse, Verständniß und Menschen Hülfe, da geht doch alles zurück, da ist doch kein Glück noch Segen; wie die Erfahrung gibt, wie der Heiden und Bibel Historien anzeigen. Ob aber unser HErr GOtt etwan den Gottlosen etliche böse Anschläge, gottlose Räthe und Gedanken läßt fortgehen, daß es scheint, sie thun, was sie wollen, sie machen es, wie sie es gelüstet, und ihre sei die Herrlichkeit und Gewalt, GOtt thue nichts dazu, so nimmt es doch ein böses Ende zuletzt, wie der 73. Psalm, V. 19., sagt: „Sie gehen unter, und nehmen ein Ende plötzlich, und mit Schrecken.“

\begin{schriftzitat}
V. 4. Ohne Wandel sind die Werke des Felsen: denn alle seine Wege sind Gerichte. Treu ist GOtt, und kein Böses an ihm, gerecht und redlich ist er.
\end{schriftzitat}

\quellseite{936}
18. Er nennt mit einem fast keinen, tröstlichen Titel GOtt einen Felsen, daß er ein recht starker, gewisser Trost ist aller, die sich auf ihn verlassen und fröhlich wagen. Alle andere Hülfe und Trost ist ein ungewißer Sandgrund. Das findet sich in der Erfahrung. Aber wer auf GOtt sich herzlich verläßt, der kann in allen Anfechtungen bestehen. Also sagt auch David im 18. Psalm, V. 3.: „Mein Fels, meine Burg, mein Hort, darauf ich traue.“ Item, Ps. 62, 8.: „Bei GOtt ist mein Heil, meine Ehre, der Fels meiner Stärke, meine Zuversicht ist auf GOtt.“

\begin{schriftzitat}
Seine Werke sind ohne Wandel und vollkommen.
\end{schriftzitat}

19. Er redet von den Werken GOttes, welche GOtt wirkt in aller Gläubigen\footnote{So Walch und die Wittenberger Ausgabe; Erlanger: allen gläubigen.} Herzen, und in allen, so GOtt recht fürchten und lieben. Denn wenn die Herzen durch den Glauben gereinigt und geheiligt werden, so ist da keine Heuchelei, sondern ist rechte vollkommene Heiligkeit und Reinigkeit, und eitel rechter heiliger Gottesdienst, also, daß das Herz GOtt recht erkennt, ihn groß achtet, herzlich fürchtet und liebt, ihm auch für so unzählige Wohlthat dankt, und in allen Nöthen anruft.

20. Und die Werke, so GOtt in uns wirkt, hält er gegen unsere eigenen Werke, da wir durch Werke, von uns selber erdichtet, oder durch unsere eigenen Träume GOtt wollen einen Gottesdienst nach unserm Gutdünkel anrichten, welchen er ihm soll gefallen lassen, aber durch welchen wir uns von GOtt zu fremden Göttern abwenden. Als ist gewesen, Möncherei, item, die Messen, und dergleichen unter dem Papstthum; das sind nicht Werke ohne Wandel, sondern heuchlerische, unvollkommene, verdammte Werke.

21. Weiter sagt er: „Alle seine Wege sind Gerichte.“ Alles Leben und Wandel der Gläubigen ist rechtschaffen, geht nach GOttes Befehl und Wort, nach den heiligen zehn Geboten, nicht in menschlichen Träumen oder Gutdünken. Also ist GOtt treu, gerecht, redlich. Das redet er alles von GOtt, auf die Weise, daß er anzeige, was GOtt im ersten Gebot fordert, denn er will, daß wir ihn für einen treuen, rechten, gütigen GOtt, der es väterlich meint, achten und halten sollen. Denn hier redet er nicht von göttlicher Natur oder Wesen, sondern redet davon, wofür man GOtt halten, was man von ihm im Herzen glauben soll; will also mit diesen Worten die falschen Götter und alle abgöttischen Greuel gerühret haben. Also sollte er sagen: Wer dem GOtt wahrlich dient, nämlich, daß er ihm (wie das erste Gebot lautet) herzlich vertraut und ihn fürchtet, der ist alle seiner Sachen gewiß, der bauet nicht auf sich selber, noch auf seine Werke, sondern auf GOttes Wort und Zusage, die nicht wanken noch fehlen kann; da ist das Herz zufrieden mit GOtt, da kann das Gewissen fröhlich und stille sein. Das ist der rechte, höchste, heiligste Gottesdienst, da ist kein Falsch, keine Lügen, noch Heuchelei an, sondern das Herz inwendig ist durch den Glauben heilig und rein. Darum sind auch den Reinen alle anderen äußerlichen Werke und Wandel heilig und rein.

22. Wiederum, an allen anderen Gottesdiensten, wie köstlich sie scheinen, die außerhalb der ersten Tafel, ohne Glauben und Gottesfurcht, mit allerlei äußerlichen Geberden zugehen, da ist eitel Ungewißheit, eitel Mühe und Arbeit, Herzeleid und Jammer der Gewissen. Da bleiben die Herzen und Gewissen allezeit in Zweifel, und können nimmer zu ruhen kommen; da bleibt inwendig das Herz voll Blindheit, voll Unglaubens; darum sind auch die äußerlichen Werke, damit man will die Seligkeit verdienen, wie heilig sie scheinen, eitel Greuel, eitel Unflat und Koth vor GOtt.

\begin{schriftzitat}
V. 5. Die verkehrte und verrückte Art hat es mit ihm verderbet, und sind nicht seine Kinder, um ihres Tadels willen.
\end{schriftzitat}

23. Es prophezeiet hier Moses, wie es nach seinem Tode gehen werde, daß sie werden leichtlich von GOttes Befehl und Geboten abfallen, die zehn Gebote verachten, nach ihrem Dünkel Gottesdienste anrichten, das erste Gebot mit Füßen treten, wie alle Heuchler thun, und eine unschlachtige, böse, ehebrecherische Art werden. Darum, so werden sie alsdann auch nicht mehr liebe Kinder noch GOttes Volk sein, ob sie wohl den Namen führeten, sondern GOttes Feinde und Hurenkinder, Ehebrecher und Ehebrecherinnen, um der geistlichen Hurerei willen, die sie treiben werden [Ezech. 16, 26. f. 23, 5. 11. ff.]. Also sagt er kurz vor diesem Gesang, im 31. Capitel, V. 29.: „Ich weiß, daß ihr es nach meinem Tode verderben werdet, und aus dem Wege treten, den ich geboten habe, so wird euch dann Unglück begegnen.“

\quellseite{937}
\begin{schriftzitat}
V. 6.
\end{schriftzitat}

Dankest du also dem HErrn, deinem GOtt, du närricht und unweises Volk? Ist er nicht dein Vater und dein HErr? Hat er dich nicht gemacht und bereitet?

24. Hier entbrennt nun ganz der hohe Prophet Mosis, so er ansieht die unaussprechliche Güte, und unzählige Wohlthat GOttes, und dagegen die große Undankbarkeit und schändliche Verachtung, und sagt: Ist das der Dank, daß euch GOtt hat zu eigenen, lieben Kindern angenommen, daß er euch mitten aus dem Tode herausgerissen, aus Egypten geführt, aus den Händen Pharaonis gewaltiglich erlöst hat, vor euren Augen das Schilfmeer lassen stehen wie Mauern, euch mit Himmelbrod gespeist in der Wüste etc.? Als sollte er sagen: Für das alles, daß GOtt gerecht, treu, gütig, und allenthalben ein gnädiger Vater ist, daß er euch vor allen andern Heiden, Völkern und Königreichen auf Erden erwählt hat, um euretwillen den König Pharaonem ersäuft, die Königreiche Canaan eingenommen; item, daß er so lange Zeit euch an nichts hat mangeln lassen, sondern euch, wie ein Vater oder Mutter ihr Kind, auf den Händen getragen, allenthalben geschützt und behütet; werdet ihr und eure Nachkommen ihn also bezahlen, mit so seinem Danke begegnen, daß ihr allerlei abgöttische, närrichte Gottesdienste werdet höher achten denn das erste Gebot GOttes, denn die erste Tafel, da er fordert, man solle ihn für einen GOtt halten, auf ihn allein bauen und trauen. Denn ihr werdet abtrünnige Hurenskinder werden, die auf ihre eigenen Werke, auf Götzendienst, auf Menschengunst, Hülfe, Trost, auf Gold und Silber werdet können hoffen und vertrauen; allein auf GOtt nicht, den werdet ihr leicht achten, als sei keine Hülfe noch Rath bei ihm, als sei an ihm gar nichts gelegen. Darum werdet ihr auch nicht Söhne des Vaters sein, sondern ein närrichtes, unweises Volk, das Hülfe suchen wird, da keine Hülfe ist. Das mag man uns [zu] dieser Zeit des Evangelii auch wohl sagen: Ist das der Dank, daß ihr vom Papstthum und unzähligen Beschwerungen erlöset seid? etc.

\begin{schriftzitat}
V. 7.
\end{schriftzitat}

Gedenke an die vorigen Tage, merke auf die Jahre der vorigen Zeiten. Frage deinen Vater, der wird dir es verkündigen, deine Aeltesten, die werden dir es sagen.

25. Da hebt er nun an, und zählt wie aus einem Register her die Wohlthat GOttes, und hebt an an der ersten, daß GOtt dieses Volk vor allen andern Heiden erwählt hat, daß GOtt so viel unzählige Wohlthat dem Volke gethan hat. Darum sagt er: „Frage deinen Vater“, das ist, frage deine Väter und deine Aeltesten. Als sollte er sagen: Wirst du Abraham, Isaak, und Jakob ansehen, so wirst du unzählige Wohlthaten finden, wie GOtt deine Väter geliebet, wie treulich er sich ihrer hat angenommen. Solche Erinnerung ist uns wohl vonnöthen. Denn der Art sind Adams Kinder, wenn eine Wohlthat vorüber ist, so ist sie bald vergessen, wie zu unsrer Zeit die Erfahrung gibt.

\begin{schriftzitat}
V. 8.
\end{schriftzitat}

Da der Allerhöchste die Völker austheilte, und ausstreuete der Menschen Kinder, da setzte er die Grenzen der Völker, nach der Zahl der Kinder Israel.

26. Die erste Wohlthat ist, daß wiewohl GOtt, der HErr, ein GOtt ist aller Heiden, aller Völker und Königreiche auf Erden, welche er Alle geschaffen hat, einem jeden Volke auch seine Grenzen und Maaß austheilt, wie im ersten Buch Mosis Cap. 11, 8. 9. ff. geschrieben steht, und wie Paulus auch Apost. 17, 26. 27. gedenkt, so hat er doch unter den allen kein Volk erwählt, da er seinen Gottesdienst hätte angerichtet, daß Gott und Hausvater er sich nennete, als das Volk Israel. Und ist wahrlich eine treffliche, große Wohlthat, Preis und Ruhm, daß die hohe Majestät sich nirgend will predigen, erkennen, finden, noch antreffen lassen, denn allein in Israel. Und daß unter so viel Königreichen und Landen nirgend soll GOttes Erkenntniß sein, denn allein in dem Volke, und daß er alle anderen Königreiche, Länder und Grenzen also abgeschieden hat, daß nicht weiter GOttes Volk irgend Leute auf Erden heißen sollen, denn allein so weit Israel wohnt, und ihr Land reicht; die andern alle sollen gottlose Heiden heißen.

\begin{schriftzitat}
V. 9.
\end{schriftzitat}

Denn des HErrn Theil ist sein Volk, Jakob ist die Schnur seines Erbes.

27. Das ist, weiter reicht GOttes Volk nicht, denn Jakob und Israel ist [Ps. 74, 1. 2. 100, 3.]. Mit den Worten will er anzeigen, wie gar liebe Kinder, wie gar ein heiliges, angenehmes, auserwähltes Volk sie vor allen andern Heiden sind.

\quellseite{938}
\begin{schriftzitat}
V. 10. Er fand ihn in der Wüste, in der dürren Einöde, da es heulet. Er führete ihn um, und gab ihm Verstand: Er behütete ihn wie seinen Augapfel.
\end{schriftzitat}

28. Das ist, er hat dem Volke auf dem Berge Sinai sein Gesetz gegeben, hat mit ihnen einen Bund gemacht, daß er will ihr GOtt sein; hat sie auch vierzig Jahre in der Wüste so wunderlich geführt, erhalten und ernährt, aus so viel Nöthen, Angst und Gefährlichkeit erlöst, ihren schändlichen Unglauben, ihre Ungeduld und Murren ihnen so oft zugute gehalten, ihre Sünde und halsstarrige Härtigkeit ihnen so oft vergeben. Hat sie wie einen Augapfel behütet; Trotz, der ihnen Leides thäte; hat sie vor allem Uebel, Jammer und Unglück durch seinen Engel schützen und bewahren lassen, also, daß sie an nichts Mangel gehabt. Das ist nun die andere, und eine fast große Wohlthat. Wer unter uns auch bedenkt, wie uns GOtt von Jugend auf so wunderlich führt, so gnädiglich ernährt, vor allem Uebel, wie einen Augapfel behütet, der wird auch finden, wie undankbar wir sind.

\begin{schriftzitat}
V. 11. Wie ein Adler aufwecket sein Nest, und über seinen Jungen schwebet. Er breitete seine Fittige aus, und nahm ihn, und trug ihn auf seinen Flügeln.
\end{schriftzitat}

29. Mit den Worten zeigt er an, wie er in der Wüste sie gebildet und getragen habe. Man sagt, und die Naturkündiger, als Plinius, Aristoteles und andere schreiben, daß mancherlei Adler sind, die beste Art aber pfleget seine Jungen mit den Flügeln zu schlagen, und wacker zu machen, daß sie fliegen lernen, auch gegen der Sonne Glanz zu halten, und welche nicht gleich in die Sonne sehen, werfen sie aus dem Nest, als eine Unart. Also will Moses sagen: GOtt, der HErr, hat dich unter seinen Flügeln gehalten, ausgeheget wie im Neste, und dich fliegen gelehrt,\footnote{Wittenberger und Erlanger: gelernt.} wie ein Adler seine Jungen; denn er hat es beides, mit Wohlthat und Strafe, versucht, dich dahin zu bringen, daß du fliegen lernest, das ist, ihm vertrauen und glauben; wie wir im achten Capitel\footnote{Es kann wohl kaum zweifelhaft sein, daß sich diese Anführung auf § 30 ff. der vorigen Schrift beziehe. Daraus würde folgen, daß unsere Auslegung später als 1529 zu setzen wäre. Diese Vermuthung wird zur Gewißheit durch das in § 50 dieses Capitels Gesagte.} gesehen haben.

\begin{schriftzitat}
V. 12. Der HErr allein leitete ihn, und war kein fremder GOtt mit ihm.
\end{schriftzitat}

30. Das ist, unter Mose ging der rechte wahre Gottesdienst im Schwange, da predigte, redete, und lehrte GOtt selber unter ihnen durch den Mosen, da hielt GOtt selber Haus, war nahe bei und unter dem Volk. Die andern Heiden ließ er ihre irrigen Wege im Finsterniß gehen. Also, jetzund, so die zehn Gebote gepredigt werden, und das Evangelium gelehrt, lehrt GOtt auch selber unter uns, grüßt uns freundlich genug, wenn wir ihm allein danken könnten. Vor ist eitel falsche Lehre und Irrthum gewesen. Dazu, so führte er auch äußerlich das Volk selbst, und gab ihnen selbst groß, stark, sicher Geleite, durch öffentliche, göttliche Zeichen, des Tages durch eine Wolke, des Nachts durch eine feurige Säule.

\begin{schriftzitat}
V. 13. Er ließ ihn hoch herfahren auf Erden, und äzte ihn mit dem Einkommen des Ackers, und ließ ihn Honig saugen aus den Felsen, und Oel aus den harten Steinen.
\end{schriftzitat}

31. Die dritte Wohlthat zeigt er da an, nämlich, daß GOtt hat um ihretwillen aus dem Lande Canaan die Heiden und ihre Könige vertrieben, und sie in einem herrlichen, köstlichen, guten Lande wohnen lassen, da sie vollauf und alles genug reichlich gehabt. Daß er aber sagt: „Er ließ sie Honig saugen aus den Felsen, und Oel aus den harten Steinen“, braucht er solcher Figuren, anzuzeigen, daß GOtt ihnen allerlei reichen Segen am Gewächse, an Nahrung und Früchten in den Landen gegeben hat; also, daß, wo auch nicht gute, feste Auen, wo gleich nicht Acker oder Wiesen gewesen, da sind dennoch Feigenbäume, Oelbäume, allerlei andere Bäume gewachsen; da ist dennoch Milch und Honig, und allerlei Nothdurft genug gewesen; so wachsen auch sonst Oelbäume gerne, wo dürr Land ist, und Bienen machen Honig, wo gleich nicht gut Land, sondern sandicht, dürre Land ist.

\begin{schriftzitat}
V. 14. Butter von den Kühen, und Milch von den Schafen, sammt dem Fette von den Lämmern. Und fette Widder, und Böcke mit fetten Nieren, und Weizen, und tränket ihn mit lauter Traubenblut.
\end{schriftzitat}

32. Da setzt er aber[mal] poetische, verblümte Worte, wie man in Liedern und Gesängen pflegt, und will anzeigen, daß sie mit allerlei Leibesnahrung
\quellseite{939}
und leiblichem Segen, an Vieh, an Früchten, mit Butter, mit Lämmern, Schafen, feisten Widdern, GOtt überreichlich versorgt hat.

33. Daraus sollen wir lernen, daß GOtt auch solche leiblichen Gaben gibt, und wenn er sie nicht gibt, so hilft keine Arbeit noch Fleiß zu Fuhrwerken, oder Schäfereien, kein fleißig Haushalten, Sorge oder Mühe. Welcher Edelmann oder Bauer gedenkt, daß er solche Gaben von GOtt hat, oder daß GOtt in der ersten Zeile der zehn Gebote, da er sagt: „Ich bin dein GOtt“, fordert, daß man alle dieselbigen Güter von ihm erbitten und warten soll, oder daß sie ohne ihn niemand haben oder behalten kann? „Traubenblut“ aber nennt er hier Wein, daß in denselbigen Ländern viel, und das mehrere Theil rother Wein wächst.

\begin{schriftzitat}
V. 15. Da er aber fett und satt ward, ward er geil. Du bist fett, und dick, und glatt worden. Und er hat den GOtt lassen fahren, der ihn gemacht hat, er hat den Fels seines Heils gering geachtet.
\end{schriftzitat}

34. Das ist, GOtt, der dich aus des Königs Pharaonis Hand erlöst hat, der dir so unzählig viel Gutes erzeigt hat, deß hast du vergessen, und eben dadurch vergessen, daß dir es wohl ging.

35. Es ist eine böse, giftige Art menschliches Herzens, daß es durch göttliche Wohlthat, durch Glück und gute Tage ärger wird, und je mehr GOtt gnädet, und je größere, trefflichere Gaben ihm GOtt gibt, sie sind leiblich oder geistlich, je weniger es sich gegen GOtt demüthigt, je weniger sich’s fürchtet, je sicherer es wird, je eher es GOttes und seines Wortes vergißt, je eher es auf sich selbst trauet und bauet, und will selbst GOtt sein. Darum warnt Moses in diesem Buche am 6. Capitel, V. 10–12., so treulich, und sagt: „Wenn du nun wirst haben große und feine Städte, die du nicht gebauet hast, und Häuser alles Gutes voll, die du nicht gefüllet hast, und ausgehauene Brunnen, die du nicht ausgehauen hast, und Weinberge und Oelberge, die du nicht gepflanzet hast, daß du kannst essen, und satt werden, so hüte dich, daß du nicht des HErrn vergessest, der dich aus Egyptenland geführet hat“ etc. Item, am 8. Capitel, V. 7. 8. 10. 11.: „Der HErr, dein GOtt, führet dich in ein gut Land, ein Land, da Bäche, Brunnen und Tiefen innen sind, die an den Bergen und in den Auen fließen; ein Land, da Weizen, Gerste, Weinstöcke, Feigenbäume und Granatäpfel innen sind; ein Land, da Oelbäume und Honig innen ist etc. So hüte dich, daß du des HErrn, deines GOttes, nicht vergessest.“ Und also haben auch die Heiden gesagt, es sei leichter, Unglück zu tragen, denn Glücks und guter Tage nicht mißbrauchen; und zeugen alle, es sei noch nie so ein weiser Mann gefunden, der Glücks nicht sei hoffärtiger worden. Ist es aber nicht eine schändliche Plage, daß GOtt für alle seine unzähligen Wohlthaten deß von uns soll gewärtig sein? wie wir zu dieser Zeit auch sehen.

\begin{schriftzitat}
V. 16. Und hat ihn zu Eifer gereizet durch Fremde, durch die Greuel hat er ihn erzürnet.
\end{schriftzitat}

36. Dies haben sie, leider, allzustark erfüllt. Denn, so sie sollten für so unzählige leibliche und geistliche Wohlthat dankbar sein, so führen sie zu, fielen von GOtt ab zu fremden Göttern, sündigten ohne Unterlaß mit allerlei Greuel und Abgötterei wider die erste Tafel, reizten GOtt, den HErrn, zu Eifer und Zorn, bezahlten ihn für seine Wohlthat mit eitel Ungehorsam und Undank; wie denn die Historien der Richter, auch die Bücher der Könige anzeigen. Denn also, seine Früchtlein sind Adams Kinder: Hält sie GOtt spärlich, oder läßt sie Armut leiden, so murren sie wider ihn, wie die Israeliten auch thäten, da sie sprachen, als ihnen Wasser mangelte: „Ist denn GOtt unter uns, oder nicht?“ [2 Mos. 17, 7.] Gibt er denn genug, so werden sie sicher und stolz, verachten und vergessen GOttes.

37. Darum, an GOtt, den HErrn, sich halten, und ihm herzlich vertrauen (es gehe wohl oder übel), ist die höchste Heiligkeit, die höchste Kunst und Weisheit, der edelste Gottesdienst auf Erden. Und wenn die Heuchler gleich lange mit andern Gottesdiensten sind umgegangen, so kommt doch immer die erste Zeile in den zehn Geboten, und sagt: „Ich bin dein GOtt“, du Heuchler, du heiliger Papist, du heiliger Papist, da hast du nie einen Tüttel davon geglaubt.

\begin{schriftzitat}
V. 17. Sie haben den Feldteufeln geopfert, und nicht ihrem GOtt, den Göttern, die sie nicht kannten, den neuen, die neulich gekommen sind, die eure Väter nicht geehrt haben.
\end{schriftzitat}

38. Das ist, allerlei Abgötterei der Heiden, als Baal, Ashtaroth, Moloch, Chamos, Camarim etc., und richteten allerlei Gottesdienst an,
\quellseite{940}
auf den Bergen und Höhen, in den Wäldern und Hainen, von welchen nichts im Gesetz geboten war, davon ihre Väter nichts gewußt, die sie ihnen selbst erfunden und erdichtet hatten, da sie ihrem Dünkel und Gedanken folgten, darüber hernach alle rechten Propheten heftig klagen. Also haben wir mit den Klöstern und mit der Papisten Messe auch gethan.

\begin{schriftzitat}
V. 18. Deinen Fels, der dich geboren hat, hast du aus der Acht gelassen, und hast vergessen GOttes, der dich gemacht hat.
\end{schriftzitat}

39. Diese Worte redet Moses aus großem, mächtigem Ernst, und aus einem Herzen, das entbrannt ist in göttlichem Eifer, GOttes Ehre wider des Teufels Lügen zu erretten. Also\footnote{Erlanger: also.} sollte er sagen: Ach! es möchte mir mein Herz brechen vor Jammer, daß ich weiß, und jeztund gleich wie vor Augen sehe, wie ihr den wahren lebendigen GOtt, den großen, gewissen Trost, die allerfesteste Burg und stärksten Fels in allen Anfechtungen und Nöthen, ohne welchen ihr nicht einen Augenblick vor dem Teufel leben oder bleiben könntet, so leicht verlassen, und so schändlich verachten werdet, welcher euch doch allein, wie eine Mutter ihr Kind, geboren und aufgezogen hat. Ja, daß ihr des herzlichen, lieben Vaters (der euch so väterlich meint, so gnädiglich vor dem Teufel, allem Uebel behütet, und so reichlich mit allerlei Gnaden und Gütern überschüttet) so schändlich vergessen werdet, und dagegen den heillosen, todten Götzen anhangen, und allerlei Abgötterei und Abgötter erdenken, die euch weder rathen noch helfen, trösten, schützen, noch erretten können: das ist Klagens werth.

\begin{schriftzitat}
V. 19. Und da es der HErr sahe, ward er beweget vor Zorn über seine Söhne und seine Töchter.
\end{schriftzitat}

40. Da zeigt er an, daß der Zorn soll desto geschwinder und heftiger sein, darum, daß so liebe Söhne und Töchter des treuen Vaters vergessen.

\begin{schriftzitat}
V. 20. Und er sprach: Ich will mein Antlitz vor ihnen verbergen, will sehen, wie es zuletzt mit ihnen hinaus will. Denn es ist eine verkehrte Art, es sind Kinder, da kein Glaube inne ist.
\end{schriftzitat}

41. Daß er sagt: „Ich will mein Antlitz vor ihnen verbergen“, redet er in der Person GOttes; das ist, ich will sie verblenden, also, daß sie mit sehenden Augen nicht sehen sollen etc. [Jes. 6, 10.], daß sie sollen in eitel Blindheit und Irrthum, ohne GOtt und GOttes Erkenntniß sterben und verderben, GOtt mit ihren Werken und Gottesdiensten suchen, und doch nicht finden. Denn GOttes Antlitz sehen, heißt in der Schrift: GOtt erkennen, seinen Willen und Werke sehen und merken. Also sagt der 80. Psalm, V. 20.: „HErr, zeige uns dein Antlitz, so ist uns geholfen.“ Wiederum, denjenigen ist GOttes Antlitz verborgen, die verblendet sind, daß sie weder GOtt noch sein Werk erkennen.

42. „Kinder“ aber, „da kein Glaube inne ist“, nennt er solche Heuchler, die heilig und fromm mit allerlei äußerlichen Werken sich stellen, und doch von dem höchsten, heiligsten Gottesdienst (welchen das erste Gebot fordert) gar nicht wissen noch halten wollen; wähnen, sie sind die lieben Kinder, so sie vermaledeite Bastarte sind, und GOttes ärgste Feinde.

\begin{schriftzitat}
V. 21. Sie haben mich gereizet an dem, das nicht GOtt ist, mit ihrer Eitelkeit haben sie mich erzürnet. Und ich will sie wieder reizen an dem, das nicht ein Volk ist, an einem närrichten Volke will ich sie erzürnen.
\end{schriftzitat}

43. Das ist, wie sie in der Wüste haben einen andern GOtt angenommen, das mir nicht zu leiden war, also will ich auch wiederum sie verwerfen, und ein ander Volk annehmen, das ihnen auch soll wehe genug thun. Wie denn auch geschehen ist, wie Paulus anzeigt zu den Römern am 11. Capitel, V. 11. ff., da die Heiden sind angenommen durch das Evangelium. Darum sind auch auf diesen Tag die Juden über uns aufs bitterste und grimmig zornig, daß wir sagen, sie sind nimmer GOttes Volk; sondern wir sind GOttes Volk, wie dieser Vers meldet.

\begin{schriftzitat}
V. 22. Denn das Feuer ist angegangen durch meinen Zorn, und wird brennen bis in die unterste Hölle, und wird verzehren das Land mit seinen Gewächs, und wird anzünden die Grundfeste der Berge.
\end{schriftzitat}

44. Feuer nennt er hier nicht allein ein leiblich Feuer, sondern die ganze, schreckliche, greuliche Verwüstung, die durchs Schwert und Feuer geschehen, da der König zu Babylon, und endlich
\quellseite{941}
die Römer Jerusalem so jämmerlich zerstört, vertilgt und verwüstet haben, daß im ganzen Lande, und sonderlich zu Jerusalem, kein Stein auf dem andern geblieben, daß alle köstlichen Gebäude zum Steinhaufen und Asche gemacht [Jer. 52, 13. ff. Luc. 19, 44.]. Und das zeigt er an, da er sagt, „und wird brennen bis in die unterste Hölle“, das ist, im ganzen Lande wird vom Gebäu oder Gewächs nichts übrig sein noch bleiben, sondern auch der bloße Boden und die Berge werden brennen, als wollte das Feuer durch die Erde in die unterste Hölle durchreißen; denn also geht es, wenn Verwüstung angeht.

\begin{schriftzitat}
V. 23. Ich will [alles] Unglück über sie häufen, ich will alle meine Pfeile an ihnen verschießen.
\end{schriftzitat}

45. Dasjenige Unglück ist beschrieben, und wie in einem Register erzählt im fünften Buch Mosis am 28. Capitel, V. 16. ff., und Moses selber beschreibt im nächsten Verse auch die Pfeile, und dasselbe Unglück.

\begin{schriftzitat}
V. 24. Vor Hunger sollen sie verschmachten, und verzehret werden vom Fieber und von bittern Seuchen. Ich will der Thiere Zähne unter sie schicken, und wüthende Schlangen.
\end{schriftzitat}

46. Da nennt er drei oder vier Pfeile, Hunger, Fieber, Pestilenz, wilde Thiere; und im nächsten Verse setzt er Krieg und Blutvergießen; und diese Plagen sind hernach oft angezogen in den Propheten. Dazu sagt er „wüthende Schlangen“. Das sind die rechten Straußbesen und scharfen Ruthen, damit unser HErr GOtt gar manchen großen, sichern, stolzen Bogen\footnote{\hebraeisch{זֵד} ein Uebermüthiger, Frevler.} und Verächter hat hingerichtet. Und wenn er anfähet, so geht es gar greulich zu. Da hilft denn nicht, daß Weib und Kind elendiglich weinen und heulen, klagen und schreien. Das wollen die Bürger und Bauern (so jemand Gottes Wort, Prediger und Pfarrherren verachten, und sagen: Die Pfaffen können nichts, denn die zehn Gebote predigen) erfahren.

\begin{schriftzitat}
V. 25. Von außen wird sie das Schwert berauben, und in den Kammern das Schrecken. Beide, Jünglinge und Jungfrauen, die Säuglinge mit dem grauen Manne.
\end{schriftzitat}

47. Das ist, Krieg und Unglück wird über sie kommen, und wird eitel Kleinmuth, Schrecken und bestürztes, erschrockenes Gemüth bei ihnen sein.

\begin{schriftzitat}
V. 26. Ich will sagen: Wo sind sie? Ich werde ihr Gedächtnis aufheben unter den Menschen.
\end{schriftzitat}

48. Das ist, ich will ihr Königreich und Priesterthum zu Boden stoßen, zu Grunde verwüsten [Hab. 3, 6. ff.], die oberste Stadt Jerusalem, Tempel und Ceremonien, alles in einen Haufen stoßen; daß, wo zuvor königlich Regiment, allerlei seine Ordnung, Opfer und herrliche Feste gewesen, ganz soll wüste und stille werden, daß jedermann mag sagen: Wo ist nun der Juden Königreich? Wo ist ihr Priesterthum? Wo sind ihre herrlichen Feste und Opfer? Das ist alles wahr worden, wie Jeremias, der Prophet, jämmerlich beweinet und beklagt [Klagl. 1, 12.]. So gibt es auch die Erfahrung auf den heutigen Tag, daß, obwohl noch eine Handvoll Juden vorhanden, daß doch ihr Regiment und alle ihr Gottesdienst zu Grunde zerrissen sind, daß sie kein eigen Königreich noch Priesterthum haben, wie andere Völker.

\begin{quote}
Und ihr Gedächtnis ist aufgehoben unter den Menschen.
\end{quote}

49. Das ist, sie sind allenthalben verachtete, arme Leute, die keinen Namen noch Ansehen mehr haben, daß sie etwas wären oder vermöchten, sondern sind eitel Fußhüter gegen dem, daß sie zuvor oben geschwebet und das Haupt gewesen sind.

50. Daß dieses nun GOtt vermöge und thun könne, nämlich, ein ganz, groß, herrlich Wesen, als eines ganzen Königreiches Pracht und Namen, öde, wüste, und in Eil stille zu machen, deß rühmt er sich wohl im ersten Gebot, da er sagt: „Ich bin GOtt“ etc., aber man glaubt es nicht. Der Pabst und die Papisten hätten vor fünfzehn oder sechzehn Jahren auch nicht gemeint,\footnote{Vergleiche die Anmerkung zu § 29 dieses Capitels. Durch diese Aussage wird gewiß, daß unsere Auslegung in das Jahr 1530 oder später zu setzen sei.} daß so viel großer Stifte und Klöster, da ihre Gottesdienste so herrlich gingen, sollten wüste werden um ihrer Greuel willen. Aber GOttes Gericht und Werk stehen vor Augen.

\begin{schriftzitat}
V. 27. Wenn ich nicht den Zorn der Feinde scheuete, daß nicht ihre Feinde stolz würden, und möchten sagen: Unsere Hand ist hoch, und der HErr hat nicht solches alles gethan.
\end{schriftzitat}

\quellseite{942}
51. Das redet er darum, denn er hat oft seinen Zorn aufgezogen, und hat auch die Gottlosen erlöst und von den Feinden errettet. Das hat er aber gethan um seines Namens willen, daß sein Name nicht gelästert würde von den umliegenden Heiden. Darum hebt er mit diesem Vers auf der Juden Ruhm und Stolz, da sie vielleicht möchten sagen: Wir sind nicht so böse gewesen, wie Moses singt, denn GOtt hat uns oft durch Wunderzeichen erlöst. Denselben antwortet er, und sagt: Es ist nicht um eures Verdienstes willen geschehen; denn ihr hättet die obgedachte Strafe wohl allezeit verdient; aber so ich euch hatte zu einem Volke angenommen, so habe ich meines Namens müssen verschonen, ob ihr wohl verzweifelte, ungläubige Buben wäret. Aber ich will dennoch nicht ewig so zusehen noch schonen, sondern es soll zuletzt euer Ende kommen, wie ihr verdient habt.

\begin{schriftzitat}
V. 28. Denn es ist ein Volk, da kein Rath innen ist, und ist kein Verstand in ihnen.
\end{schriftzitat}

52. Das ist, sie kennen weder GOtt, noch seine Werke; sie rühmen, ich sei ihr GOtt, und kennen mich nicht, gehen in ihrer Blindheit sicher dahin mit ihrer Heuchelei, und treiben ohne Unterlaß allerlei Abgötterei wider das erste, höchste Gebot, sind vermessen und sicher, wollen wähnen, sie sind liebe Kinder, und ihre Sache stehe wohl bei mir, so doch mein grimmiger Zorn bald soll über sie ausgegossen werden. Darum denken sie nicht einmal daran, daß sie sich bekehren, oder bessern.

\begin{schriftzitat}
V. 29. O daß sie weise wären, und vernähmen solches, daß sie verstünden, was ihnen hernach begegnen wird.
\end{schriftzitat}

53. Als sollte er sagen: Ich habe so lange ihnen gepredigt, jetzund habe ich sie genug vermahnt, daß sie (wie das erste Gebot lautet) den einigen GOtt allein fürchten und lieben sollten, und auf ihn vertrauen, und von ihm nicht weichen. Die Propheten nach dieser Zeit werden auch nichts Anderes treiben, denn daß die Furcht GOttes und der Glaube gegen GOtt der höchste Gottesdienst sei. Wollte GOtt, wenn solche Predigten gehen, sie würden dann weise, verachteten es nicht, und erkenneten, daß ihr Tod und Leben, ihr Gedeihen und Verderben, ihre Seligkeit und Verdammniß ihnen daran gelegen ist.

54. Aber ich weiß, sie werden in Sicherheit und Abgötterei gerathen; und wie sie mir haben viel Herzeleid gemacht in der Wüste, also werden sie die künftigen Propheten auch in den Wind predigen lassen, nichtsdestoweniger ihren Abgöttereien nachlaufen, bis so lange Sennacherib und Salmanasser, item, Nebucadnezar, der König zu Babel, kommen, und endlich die Römer, und lehren\footnote{Mittenberger und Erlanger: lernen.} sie die zehn Gebote verstehen, nämlich, daß GOtt ist ein eifriger, starker GOtt, der endlich Unglauben und gottlos Wesen nicht ungestraft läßt. Also sagt Salomon auch in Sprüchen am 1. Capitel, V. 20. 24. ff.: „GOttes Weisheit ließ sich hören auf den Gassen. Sie redt ihre Hand aus, und niemand achtet darauf. So will ich auch lachen (sagt GOtt) in eurem Unfall, und euer spotten, wenn da kommt, das ihr fürchtet; wenn über euch kommt, wie ein Sturm, das ihr fürchtet, und euer Unfall als ein Wetter, wenn über euch Angst und Noth kommt.“ Also, da der HErr Christus die Stadt Jerusalem mit Erbarmen ansah, und über sie weinte [Luc. 19, 41.]. Da er selber, hernach auch die Apostel, das jüdische Volk zur Besserung vermahnten, und sie vor ihrem großen Unfall warnten, waren gar wenig Weise, die es zu Herzen nahmen, hatte niemand Achtung darauf, was nach\footnote{Erlanger: noch.} ihrem Predigen oder herzlichen Klagen und Weinen folgen würde; aber hernach eine kurze Zeit erfuhren sie es allzustark.

\begin{schriftzitat}
V. 30. Wie gehet es zu, daß Einer wird ihrer tausend jagen, und zwei werden zehntausend flüchtig machen? Ist es nicht also, daß sie ihr Fels verkauft hat, und der HErr hat sie übergeben?
\end{schriftzitat}

55. Das ist: Wenn es einmal zur Strafe kommen wird, wird ihnen Herz und Muth entfallen; so werden sie nicht allein aus der Propheten Worte hören, sondern mit dem Werk erfahren, daß es aus mit ihnen ist, daß sie GOtt verlassen hat. Und also ist es auch gegangen, da der König von Babel, wie ein schrecklich Wetter sie überfiel; da die Römer Jerusalem in Grund verstörten, daß Ein Heide tausend Juden schreckte, und in die Flucht schlug, und daß auf der Juden Seite nichts denn Zagen, Schrecken, Zittern, Beben und Verzweifeln war, so doch zuvor, da GOtt noch bei ihnen war, Ein Jude tausend Heiden schreckte, und aller gläubiger Segen auf ihrer Seite war [3 Mos. 26, 7. 8.]. Also, [zu]
\quellseite{943}
dieser Zeit, wird man nicht eher glauben, was für eine starke Macht GOttes Beistand ist, bis daß er die Hand abthut, und uns versuchen läßt, was Stärke, Vermögen, Herz oder Muth wir von uns selber haben.

\begin{schriftzitat}
V. 31. Denn unser Fels ist nicht wie ihr Fels, deß sind unsere Feinde selbst Richter.
\end{schriftzitat}

56. Das ist, da euer GOtt, euer großer, gewisser, starker Schutz, bei euch hielt, da hattet ihr gut thun, da war ein König der Heiden euch wie eine Fliege, da galt es gleich, mit großer oder kleiner Rüstung den Feinden begegnen, denn GOtt war bei euch, der nicht allein Roß und Mann, sondern auch alle Herzen und Gedanken der Feinde in der Hand hat.

\begin{quote}
Denn unser Fels ist nicht wie ihr Fels.
\end{quote}

57. Das ist, wir haben wohl einen andern HErrn und GOtt denn die Heiden. Die Heiden haben silberne und güldene Götzen, haben viel Reichthümer und Macht; aber unser GOtt hat es alles in der Hand, und kann niemand ein Schwert stürzen, ja, niemand euch sauer ansehen, oder einen bösen Gedanken haben, er wolle es denn.

58. Das alles haben die Heiden selbst müssen bekennen, daß dem Volk Israel mit Gewalt oder Macht niemand konnte anhaben oder abbrechen, so lange ihr GOtt bei ihnen hielte. Darum gab auch Bileam dem König Moab einen Rath, daß er das Volk erst sollte zu Sünden reizen, daß ihr GOtt von ihnen ließe, und er also obsiegen möchte [4 Mos. 31, 16.]. So ist nun Bileam einer von den Zeugen, die da zeugen müssen, daß Israels GOtt ein starker Fels und unüberwindlicher GOtt sei. So sind die Egypter auch Zeugen, die da sagten im 2. Buch Mosis am 14. Capitel, V. 25.: „Lasset uns [von] Israel fliehen, denn GOtt streitet für sie.“ So sind auch Zeugen die Philister, im ersten Buch Samuels am 5. Capitel, V. 7. Darum haben seine Götter so hart bei ihrem Volke gehalten, oder ihre Gewalt bewiesen, als der wahre, lebendige GOtt in Israel. Also sagt er hier oben im 3. Capitel, V. 22.: „Fürchtet euch nicht vor ihnen, denn der HErr, euer GOtt, streitet für euch.“

\begin{schriftzitat}
V. 32. Denn ihr Weinstock ist des Weinstocks zu Sodom, und von dem Acker Gomorra. Ihre Beere ist eine Gallenbeere, sie haben bittere Trauben.
\end{schriftzitat}

59. Das ist, sie sollten wohl ein feiner, lustiger Weinberg sein, den GOtt selbst gepflanzt hat, darinnen viel fruchtbare Weinstöcke nach aller Lust stünden [Hos. 10, 1.], sie sollten wohl ein fein gottesfürchtig, dankbar Volk sein, da viel, viel gläubiger, frommer Kinder unter wären, die GOttes Wort von Herzen lieb hätten, andere Leute auch herzu brächten, und viel Frucht trügen: so sind es eitel Disteln und Dornen, eitel wilde Trauben, eitel Mißwachs, faule, verdorbene Beeren, und eitel solche feine Früchtlein, wie Sodoma und Gomorra. Es sollte eitel köstlicher, reiner, kräftiger Most und guter Wein unter ihnen gefunden werden: so ist es eitel Lauer,\footnote{Im Original und in der Wittenberger richtig: „Lauer“, was Walch und die Erlanger in „sauer“ verwandelt haben. „Lauer“ ist schlechter Wein. Siehe Bd. XX, 2396 unserer Ausgabe s. v. Lauer.} ja, eitel Gift und Galle, allerlei falsche Lehre, und selbst erdichtete, abgöttische Gottesdienste. Also haben die andern Propheten solch Gleichniß aus Moise genommen; als, da Jesaias am 5. Capitel, V. 2., sagt: „Er wartete, daß sein Weinberg Trauben brächte, aber er brachte wilde Trauben.“ Und der 80. Psalm, V. 14., vergleicht auch das ganze Volk einem Weinberge, welchen die wilden Säue zerwühlen etc. Diese heftigen, harten Worte redet er aus ganz hitzigem Eifer und großem Ernst, wider die falsche Lehre und falschen Gottesdienst.

\begin{schriftzitat}
V. 33. Ihr Wein ist Drachen Grimm, und wüthiger Ottern Galle.
\end{schriftzitat}

60. Diesen feinen Titel (da er es abermal aus großem, heftigen Zorn und Eifer Drachen Grimm, Gift und Galle nennt) gibt er ihrer falschen Lehre, dadurch allerlei greulicher Irrthum ward eingeführt, und die rechte Lehre, vom Glauben, der wahren Propheten unterdrückt. Denn es ist kein schädlicher und ärger Gift unter der Sonne, denn solche Lehre, welche mördlichen und\footnote{So die Wittenberger. Erlanger: mördlich unsäglichen.} unsäglichen Schaden thut, und die Leute ohne Unterlaß weiter von GOtt führt, zu eitel Greuel und eitel Gotteslästerung.

61. Solche Gift und Galle war der Gottesdienst Baal, zu des Propheten Elias Zeiten [1 Kön. 18, 26.]. Solche Gift und Galle waren auch die Kälber Samaria, und alle Gottesdienste auf den Höhen etc. [1 Kön. 12, 28. 32.]. Noch hing das Volk, die Könige und Fürsten denselbigen
\quellseite{944}
an, hörten die falschen Propheten, verfolgten und corrigirten die frommen Propheten [2 Chron. 18, 17. ff.]. Solcher Gift sind unter dem Papstthum so viel unzählige, gottlose Lehren gewesen: die Lehre von der Messe und vom opere operato, vom Fegfeuer, von Klostergelübden, von dem ehelichen Stande etc., welche schrecklichen und unaussprechlichen Schaden gethan; noch wird das alles durch den Teufel und seine Apostel eitel heilig Opfer, heilige, göttliche Andacht und Gottesdienst gerühmt, so es hier der Heilige Geist durch Mosen schädliche Gift und Galle nennt; und Paulus der Apostel sagt auch [1 Tim. 4, 1. 2 Tim. 2, 26.], solche falsche Lehre sei teuflische Lehre, und fresse um sich wie der Krebs.

\begin{schriftzitat}
V. 34. Ist solches nicht bei mir verborgen, und versiegelt in meinen Schätzen?
\end{schriftzitat}

62. Das ist, haltet es nicht dafür, daß ich ein GOtt sei wie Baal, der nicht sehe noch höre, vor dem ihr euch könnet endlich decken oder verbergen. Ich weiß und sehe allen Greuel, allen euren Wandel und Werke, die ihr treibt; allein um meines Namens willen schone ich noch, aber geborget soll es euch sein, und nicht geschenkt. Ob ihr wohl nicht seht aus Blindheit und Sicherheit, was für eine Strafe über euch im Himmel beschlossen ist, und wie das erste Gebot soll und muß wahr bleiben, so habe ich doch alle Register eurer Greuel, eurer Undankbarkeit, eurer Verachtung bei mir verborgen, und darf es ihm keiner in Sinn nehmen, daß ein Verächter oder Gottloser sollte ungestraft bleiben.

63. Dieses ist wider Sicherheit geredet, daß sie nicht denken, dieweil ihnen GOtt Geld, Güter, Acker und Wiesen gibt, daß sie darum mit GOtt wohl stehen, und alle Sache gut sei; denn leibliche Güter gibt er auch den Heiden; sondern, daß sie bedenken, was sie wohl für Strafe verdient haben, wie das erste Gebot soll wahr bleiben, als es nicht kann falsch sein, nachdem GOtt, dem HErrn, seine göttliche Majestät, Ehre, Preis, und ganze Gottheit daran gelegen ist. Denn es darf niemand denken, daß er darum einen gnädigen GOtt habe, wenn er ihm leibliche Güter gibt, Haus, Hof, Weib, Kind, Acker, Wiesen; denn da sagt der 37. Psalm, V. 20. 35., und alle Schrift, daß die Gottlosen in der Welt am reichsten und glücklichsten sind; sondern eben wenn solcher zeitlicher, leiblicher Segen vorhanden ist, so sehe man am ersten darauf, was GOtt im heimlichen Register hat.

64. Aber das kann kein Gottloser thun, sondern gehen in eitel Blindheit und Sicherheit dahin. Also jemand Adel, Bauern und Bürger etc., daß sie vieler Bürden entledigt sind, des Ablasses, des Bauwes, der Geldmessen, der Stationirer, der Bettelmönche etc., durch welches ihnen, über das zeitliche, Heil der Gewissen, jährlich merklich an ihrer Nahrung zugeht, dadurch sie manchen Gulden inne behalten; meinen, so GOtt ihnen die leibliche Nahrung mehrt, sie haben einen gnädigen GOtt; gehen hin in eitel Sicherheit, verachten nun auch aufs äußerste GOttes Wort und die Prediger, die es treulich meinen; wissen aber nicht, was GOtt im Register hat, und daß er solcher Sünde schreckliche Strafe spart, bis zu seiner Zeit.

\begin{schriftzitat}
V. 35. Die Rache ist mein, ich will vergelten. Zu seiner Zeit soll ihr Fuß gleiten; denn die Zeit ihres Unfalls ist nahe, und ihr Künftiges eilet herzu.
\end{schriftzitat}

65. In diesem Worte, das GOtt hier durch Mosen sagt, „die Rache ist mein“, wird abermal gerühmt von dem ersten Gebot, daß es soll wahr bleiben, und kein Gottloser soll der Rache oder Strafe entfliehen. Also redet diesen Worten Mosis nach der 94. Psalm, V. 1–3.: „Herr Gott, des die Rache ist, Gott, des die Rache ist, erscheine. Erhebe dich, du Richter auf Erden, wie lange sollen die Gottlosen trotzen?“ „Mein (sagt er) ist die Rache“; als sollte er sagen: Ob ich eine Zeitlang schone und die Strafe aufziehe um meines Namens willen, so soll euer Greuel doch nicht ungestraft bleiben; ob ihr wohl jetzund in einem guten Lande wohnt, Aecker und Wiesen, Weingärten und allerlei Früchte reichlich habt, und also wie in einer Mummerei oder Larve einhergehet, als seiet ihr liebe Gotteskinder, so ist doch die Rache mein, und meiner Hand, wenn die einmal gehen wird, soll niemand wehren; wenn ihr am sichersten seid, so will ich mich sehen lassen, daß mein Name sei Herr.

66. Also ist ihnen auch widerfahren: Die Strafe ward aufgezogen, bis da das Endkindlein kam, an welchem die Römer das Garaus mit ihnen spielten, und sie auf einmal dahin suchten, wie sie wohl lange zuvor verdient hatten. Und also ist ihr Fuß geglitten, daß sie nun weder
\quellseite{945}
Königreich noch Priestertum haben. Und also siehst du, daß Mosis prophezeit, daß einmal das Volk mit Königreich und Priestertum soll ein Ende haben. Die guten Gesellen, so jetzund das Evangelium so schändlich verachten und übel mißbrauchen, werden auch erfahren, daß GOttes die Rache ist.

\begin{schriftzitat}
V. 36. Denn der HErr wird sein Volk richten, und über seine Knechte wird er sich reuen lassen. Denn er wird ansehen, daß ihre Kraft dahin ist, und aus ist, auch mit dem, der verschlossen und übrig war.
\end{schriftzitat}

67. Das ist, in dieser Sache, da es seine Ehre belangt, wird er auch seines eigenen Volks nicht schonen, und wird nicht ansehen den herrlichen Namen und Ruhm, daß sie Israel, daß sie GOttes Volk und GOttes Erbe heißen, daß sie liebe Kinder heißen, sondern des ersten Gebots Urtheil, das da heißt: „Er sucht daheim die Missethat der Väter über die Kinder“ etc. wird sie treffen. Und da soll nichts für helfen. Gnädig wird er aber sein, wie der andere Theil des ersten Gebots lautet, und Barmherzigkeit thun, ohne alles Ansehen der Person, denjenigen, die ihn von Herzen fürchten und lieben, und seine Gebote halten. Und hier prophezeit Moses, daß der übrige kleine Haufe Israel, so sich zu Christo bekehren wird, soll erhalten werden. Sonst heißt es also [V. 36.]: „Ihre Kraft soll ganz und gar aus sein, auch mit dem, der heimlich verschlossen und übrig war“, das ist, es soll vom Königreich Israel, und endlich vom ganzen Volk, wenig oder nichts überbleiben, alle sollen sie unter die Heiden in Jammer und Elend zerstreuet werden.

\begin{schriftzitat}
V. 37. Und man wird sagen: Wo sind ihre Götter? Ihr Fels, darauf sie vertrauten?
\end{schriftzitat}

68. Mit allen Gottlosen und solchen Heuchlern, welche Heiligkeit und Gottesdienst vorgeben, ohne Gottesfurcht und Glauben, ohne den höchsten Gottesdienst der ersten Tafel, läßt es GOtt endlich dahin kommen, daß ihnen ihr eigen Gewissen widerspricht und ihrer spottet, daß sie in ihren großen Aengsten denken: O wo bleiben wir nun? O wo ist nun unser GOtt? Denn, dieweil es wohl geht, sind alle Heuchler ganz sicher, meinen, sie stehen mit GOtt sehr wohl, denken nicht anders, es gehe ihnen darum so wohl, daß sie so heilig sind; wenn aber Widerwärtigkeit kommt, so werden sie verzagt, und sagt ihnen denn ihr eigen Herz und Gewissen also: Was helfen uns nun unsere guten Werke, unsere Arbeit und Mühe? [Weish. 5, 6. 7. ff.]

69. Wahr ist es, der Teufel wollte ein solches bitteres Wort den rechten Christen und Gläubigen auch gerne ins Herz reden: Wo ist nun euer GOtt? wie der 42. Psalm, V. 4., anzeigt. Aber der wahre GOtt hält fest über seinem ersten Gebote, und wenn es aufs höchste, oder zu den rechten Zügen und Nöthen kommt, bleibt er endlich nicht außen. Er sagt bald, wie still er schweigt: Hier bin ich; und je länger er verzeucht, je reicher er tröstet. Aber falscher, abgöttischer Trost, der bleibt außen, und die falschen Götter lassen ihre Anbeter schweben und sinken. Da sehen und merken sie denn erst, was sie in ihrer Sicherheit gewesen sind, was sie geglaubt haben, wie einen feinen Gott sie haben. Das meint er: „Man wird sagen: Wo sind eure Götter?“

\begin{schriftzitat}
V. 38. Von welcher Opfer sie Fett aßen, und tranken den Wein ihres Trankopfers. Lasset sie aufstehen und euch helfen, und euch schützen.
\end{schriftzitat}

70. Sie lassen sich dünken, sie seien sehr heilig, und essen nicht Götzenopfer, sondern ihre Opfer gefallen dem wahren, rechten GOtt. Also haben die Papisten gethan mit dem großen Mißbrauch der Messen. Also rühmt der Türke, welcher seinen Mahomet über und neben Christum setzt, er diene auch GOtt, und habe unter seinem Volke große Gottesdienst; aber dieweil er nicht Christum erkennt, und den wahren GOtt nicht glaubt, ist alle sein Gottesdienst ein Greuel, und man wird ihm auch sagen: Wo sind nun ihre Götter? Wo ist nun ihr Mahomet? Und nächst, da er vor Wien in Oesterreich gelegen ist,\footnote{Vergleiche in der vorigen Auslegung Cap. 7, § 54. Dort heißt es: „vor wenig Tagen“; hier: „nächst“, das heißt, in jüngst vergangener Zeit.} hat man es wohl sagen mögen, da ihm in die sechzig oder achtzig tausend Menschen an der Pestilenz gestorben sind, und da ein solch Verzagen in sein Volk kommen, daß er hat müssen abziehen, und des wahren, lebendigen GOttes und unseres Heilandes Jesu Christi Hand und Macht öffentlich sehen und erfahren.

\begin{schriftzitat}
V. 39. Sehet ihr nun, daß ich Ich bin, und ist kein GOtt neben mir. Ich kann tödten, und lebendig machen; was ich zerschlagen habe, daß
\end{schriftzitat}

\quellseite{946}
kann ich heilen, und ist niemand, der es aus meiner Hand errettet.

71. Das ist der Gottlosen und aller Heuchler Art, daß sie Heiligkeit vorgeben mit vielen äußerlichen Gottesdiensten, als achten sie groß auf GOtt, und leben doch in aller Sicherheit, verachten alle Gebote GOttes, sonderlich das erste Gebot, als sei kein GOtt, als seien die zehn Gebote ein Traum (Tit. 1, 16.). Darum, wo GOtt, der HErr, seine Gottheit soll erhalten, so muß er endlich durch Erfahrung sie solches finden und erfahren lassen, sonst sprechen sie immer, wie der Psalm sagt in der Person der Gottlosen: Sollte davon GOtt etwas wissen? [Ps. 14, 1. 4. Ps. 53, 1.] Darum muß er durch die Erfahrung beweisen, daß er GOtt sei, daß also die Verächter endlich bekennen müssen, er lebe und herrsche noch, wenn sie ihn lange verachtet haben.

\begin{schriftzitat}
V. 40. Denn ich will meine Hand in den Himmel heben, und will sagen: Ich lebe ewiglich,
\end{schriftzitat}

72. Dies ist aber eine starke Bekräftigung, da GOtt, der HErr, durch sich selbst schwört, daß die erste Tafel und das erste Gebot soll wahr bleiben. Und dieser Eid, da GOtt, so er keinen Größern über sich hat, bei sich selber schwört, wird hier ausgedrückt zu Trost den Gläubigen und wahren Israeliten, und zum Schrecken den Verächtern und Gottlosen. Denn dieser starke Eid, und solcher herrlicher Trost und Ruhm zeigt an, daß GOtt allein ein gewisses, ewiges, unendliches Leben und unendliche Macht und Gewalt hat. Alle unser Leben und Wesen, Thun und Vermögen ist ungewiß, vergänglich. Allhier wird nun GOttes ewige Macht, die niemand hindern kann, gerühmt, wie man sagt: Richten, Rühmen, Rächen; die drei R gehören GOtt allein; wir sollen von uns selbst uns nichts unterstehen, noch von uns etwas rühmen. Also sagt der 62. Psalm, V. 12.: „GOttes ist die Macht.“

\begin{schriftzitat}
V. 41. Wenn ich den Blitz meines Schwerts wetzen werde, und meine Hand das Gericht angreifen wird, so will ich mich wieder rächen an meinen Feinden, und denen, die mich hassen, vergelten.
\end{schriftzitat}

73. Das ist, das Wort, so im ersten Gebot steht, daß Ich, GOtt, dahin suche die Wissenschaft der Väter über die Kinder, in das dritte und vierte Geschlecht etc., das will ich wahr machen und behalten, nicht allein an euch, meinem Volke Israel, sondern in allen Königreichen der Heiden; und braucht eben die Worte des ersten Gebots: „Denen, die mich hassen“ etc.

74. Daß er aber sagt: „Den Blitz meines Schwerts“, will er anzeigen, daß GOttes Strafe und Rache solle groß und schrecklich sein, die da wider alle Gottlosen, wider alle Verächter göttlicher Gebote, sie sind gleich Juden oder Heiden, wenn das Stündlein kommt, hindurch reißen und dringen soll, wie ein Blitz und starker plötzlicher Donnerschlag, der nicht zu hindern, denn nicht zu wehren ist.

\begin{schriftzitat}
V. 42. Ich will meine Pfeile mit Blut trunken machen, und mein Schwert soll Fleisch fressen, über dem Blut der Erschlagenen, und über dem Gefängniß, und daß des Feindes Haupt entblößt sein wird.
\end{schriftzitat}

75. Meine Strafe und greuliche Rache will ich also machen, daß, ob ich wohl den Zorn habe aufgezogen, [es] doch aufs letzte ganz erbärmlich soll zugehen. Daß aber Moses in der Person GOttes sagt: „Meine Pfeile sollen voll Bluts sein“, „mein Schwert soll Fleisch fressen“, redet er aber[mal] auf der hebräischen Sprache Art. [Er] will aber sagen: Wenn der Zorn angeht, will ich plötzlich ihren ärgsten Feinden Raum lassen, Krieg, Jammer und Unglück über sie, wie einen starken, großen Strom, gehen lassen, und ein Blutvergießen anrichten, daß da soll kein Erbarmen noch Schonen sein.

76. Dies alles ist auch, wahrlich, mit den Juden also gegangen, da der König zu Babel, NebucadNezar, da ernach\footnote{Im Original und in der Erlanger: „da ernach“; Wittenberger: „da hernach“.} die Römer, die Stadt Jerusalem verwüsteten, daß, wie Josephus schreibt, todt blieben erstmal hunderttausend durch Hunger und Schwert. Daß er aber sagt: „Das Haupt des Feindes soll entblößt werden“, das ist, mein ungehorsam Volk, meine Feinde, sollen ihres Königreichs beraubt werden, [es] soll alles mit ihnen ein Ende haben. Wir sehen zu unserer Zeit auch öffentlich, daß GOtt durch allerlei Zeichen seinen Zorn dräuet mit dem allerschrecklichsten Feinde, dem Türken; so hat er, leider, viel christliches Bluts vergossen; aber es sind wenige, die dem göttlichen Zorne mit ernsten Gebet entgegen gingen.
